\subsection{Selective, structure-preserving editing}
\label{sec:exp_sel}
% generated by make_tab_selective.py from align/sel2/ev1 (one evaluation job for every row)
\begin{table*}[t]
\centering
\caption{\textbf{Selective, structure-preserving editing} (pre-registered; 526 clips; one front-end pass, in which the 99 two-row clips were re-run after a compositing fix, and one evaluation run for every row: all rows read together, every metric computed over all rows at once). B0 is \method without the preservation gate; each component is added to B0 separately and then all three together (B4). Collateral: change outside the old and new glyphs in the text neighbourhood. The verdict applies the pre-registered criteria; development clips are excluded in the right block. The B1 of this pass and the \method of \cref{tab:main} are the same configuration in two independent front-end passes and evaluation runs.}
\label{tab:selective}
\resizebox{\linewidth}{!}{%
\begin{tabular}{@{}lccccccccc|ccc@{}}
\toprule
 & \multicolumn{9}{c|}{All 526 clips} & \multicolumn{3}{c}{486 clips (no development clips)} \\
 & Read.\ $\uparrow$ & Fixed & Broken & Net & Appear.\ $\downarrow$ & BG PSNR & Collateral $\downarrow$ & Flicker & Verdict & Read.\ $\uparrow$ & Net & Collateral $\Delta$ vs.\ B0 \\
\midrule
B0: \method w/o preservation gate & 91.3 & 214 & 16 & 198 & 2.52 & 42.97 & 1.55 & 1.76 & -- & 91.2 & 178 & -- \\
B1: + preservation gate (= \method) & \textbf{91.6} & 214 & \textbf{14} & \textbf{200} & 2.39 & 42.97 & 1.52 & 1.79 & adopted & \textbf{91.4} & \textbf{179} & $-0.031$ [$-0.06$, $-0.01$] \\
B2: + glyph-only erasure & 83.7 & 181 & 23 & 158 & 2.43 & 43.06 & 1.46 & 1.79 & rejected & 84.4 & 145 & $-0.036$ [$-0.09$, $+0.03$] \\
B3: + carrier-fit check & 88.8 & 200 & 15 & 185 & 2.39 & 42.95 & 1.54 & 1.77 & rejected & 88.9 & 167 & $-0.004$ [$-0.01$, $-0.00$] \\
B4: all three & 82.9 & 173 & 19 & 154 & \textbf{2.21} & 43.05 & \textbf{1.43} & 1.82 & rejected & 83.3 & 140 & $-0.062$ [$-0.12$, $+0.01$] \\
\bottomrule
\end{tabular}}
\end{table*}


Inspecting the outputs of the first full evaluation, we found that the most frequent problem of \method was not that its glyphs lacked material but that it changed things it should not have: already-correct text, sign frames, and the structure around the letters. We therefore tested three components that make editing more selective, each added to B0 (\method without the preservation gate) on its own, with acceptance criteria fixed before the run (supplement): the preservation gate of \cref{sec:pres} (B1), \emph{glyph-only erasure} (B2), and a \emph{carrier-fit check} (B3).

\paragraph{Glyph-only erasure (B2).} Instead of inpainting the whole letter band, B2 erases only the old glyphs and their own glow, estimated with a mixture model of carrier, ink and halo colours, and protects sign frames, long bars and the boxes of other text. It is active on the 325 clips with an accepted instance---in 117 of them it falls back to band erasure minus the protected structures, because the surroundings themselves look like strokes (\eg neon stripes)---and its erased area is a median 0.64 of the band area. \emph{Pre-registered criterion}: collateral change decreases ($p<0.05$) and readability drops by at most one point.

\paragraph{Carrier-fit check (B3).} B3 estimates the extent of the carrier from colour, lightness and texture boundaries and declares a clip \emph{layout-infeasible} if the target, at the character density of the source, would have to shrink below 0.7 of the source height to fit; such clips are left unedited and counted as failures. \emph{Criterion}: the infeasible flags are confirmed by manual inspection.

\paragraph{Results.} \Cref{tab:selective} reports all variants from one front-end pass (with the 99 two-row clips re-run after a compositing fix that affects all variants alike). The preservation gate meets its criteria (no false skip; two fewer broken clips; readability not lower) and is adopted, which makes B1 the \method of \cref{tab:main}. Glyph-only erasure reduces collateral change on all clips ($-0.090$, 95\% CI $[-0.17,-0.02]$, $p=0.017$) but not significantly once the development clips are removed ($p=0.23$), and it costs 7.3 points of readability: the instance ink mask is not a complete glyph mask, so parts of the old strokes survive next to the new ones, and OCR boxes of the garbled text itself were sometimes protected as ``other text''. A second version that completes the old glyphs by colour still left remnants in three of nine diagnosed clips, which suggests that this component needs a learned stroke segmentation. Of the 24 clips flagged infeasible by the carrier-fit check, manual inspection confirmed 5, was undecided on 3 and rejected 16 (OCR missed letters of the source, or two-row signs were treated as one row), so B3 is rejected as well; a stricter version that only flags clips whose source stably shows a whole-word part of the target behaved correctly on the diagnosed clips but has not been run at scale. All three together (B4) inherit B2's readability loss. A blind study comparing B1 with \steer (60 clips) and B2 with B0 (28 clips) on naturalness and on damage to non-target content is in progress \pending{(results pending)}.
